% Source: https://artofproblemsolving.com/community/c3h1758669_a_quick_exercise_on_descartes_rule_of_signs
% Stable identity: aops-topic-1758669-q01
Let $p(x)$ be the polynomial. Descartes' rule says the number of positive roots (with multiplicity) equals the number of sign changes in the coefficient sequence, or is less than it by an even number; negative roots are counted the same way using $p(-x)$.

\textbf{Positive roots.} The signs of the coefficients of $p(x)$ are
\[
+\,+\,-\,+\,-\,-\,+\,- ,
\]
with sign changes $+\to-$, $-\to+$, $+\to-$, $-\to+$, $+\to-$: five changes. So $p$ has $5$, $3$ or $1$ positive roots.

\textbf{Negative roots.} $p(-x)=-27x^7+51x^6+300x^5+x^4+49x^3-52x^2-178x-37$, with signs
\[
-\,+\,+\,+\,+\,-\,-\,- ,
\]
which has two sign changes. So $p$ has $2$ or $0$ negative roots.

Answer: $\boxed{5,\ 3\text{ or }1\text{ positive roots, and }2\text{ or }0\text{ negative roots}}$.

(Evaluating, $p(-5)<0$, $p(-1)=107>0$ and $p(0)=-37<0$, so there are in fact exactly $2$ negative roots; and $p(0)<0$, $p\!\left(\dfrac{1}{2}\right)>0$, $p(1)=-181<0$, $p(3)>0$ give at least $3$ positive roots. Numerically the real roots are about $-4.42$, $-0.91$, $0.23$, $0.71$ and $2.55$.)
